% Generated by scripts/analysis/revision_reference_tables.py; do not hand-edit numbers.
\section{参考值、数值稳定性与比较敏感性}
\label{sec:reference-precision}
本节直接整理冻结参考和主要实验的保存统计量；参考偏移分析与后验区间摘要是收到本轮意见后的伴随分析，不改变原参考、可用集合或实验协议。
\subsection{W1的逐函数参考}
在驻点附近采用固定坐标$\theta=m+Cz$，其中
\[
m=(0.6059593596,-0.6218819313)^\mathsf T,\qquad
C=\begin{pmatrix}0.0603102245&0\\-0.0768479318&0.0598848727\end{pmatrix}.
\]
驻点由解析NumPy梯度求得，最大梯度绝对值为$8.68\times10^{-14}$；$CC^\mathsf T$为该点观测信息的逆。
在$[-R,R]^2$内作张量Gauss--Legendre求积，$R\in\{6,8,10,12\}$，每轴阶数$n\in\{32,64,96\}$，保留全部12个结果。事件$\beta>0$另按边界分段积分。使用预先选定的$(R,n)=(12,96)$，没有按与MCMC的接近程度选择参考。

\begin{table}[htbp]\centering\small
\begin{tabular}{lrrrrr}\toprule
函数 & 求积参考 & 加密变化 & 扩域变化 & MCMC减求积 & MCMC MCSE\\\midrule
$\alpha$ & $0.606577112$ & $4.44\times10^{-16}$ & $2.22\times10^{-16}$ & $1.33\times10^{-4}$ & $6.33\times10^{-4}$\\
$\beta$ & $-0.622983250$ & $2.22\times10^{-16}$ & $6.66\times10^{-16}$ & $8.40\times10^{-4}$ & $0.00104$\\
$\alpha^2$ & $0.371577055$ & $2.22\times10^{-16}$ & $5.55\times10^{-16}$ & $1.36\times10^{-4}$ & $7.68\times10^{-4}$\\
$\beta^2$ & $0.397618708$ & $2.22\times10^{-16}$ & $7.77\times10^{-16}$ & $-0.00109$ & $0.00135$\\
$p(0)$ & $0.647037397$ & $6.66\times10^{-16}$ & $0$ & $3.13\times10^{-5}$ & $1.44\times10^{-4}$\\
$p(100)$ & $0.495902581$ & $2.22\times10^{-16}$ & $2.22\times10^{-16}$ & $2.43\times10^{-4}$ & $1.58\times10^{-4}$\\
$\ind(\beta>0)$ & $4.767209\times10^{-11}$ & $1.77\times10^{-24}$ & $4.76\times10^{-24}$ & $-4.77\times10^{-11}$ & 未定\\
\bottomrule\end{tabular}
\caption{W1参考精度摘要。加密变化为$R=12$时64到96阶之差；扩域变化为96阶时$R=10$到12之差。后两列来自独立的既有Stan参考；MCSE取posterior、链间均值和批均值三种估计的可用最大值。上述变化是数值稳定性指标，不是认证误差界。}
\label{tab:reference-wells}
\end{table}
独立MCMC参考来自posteriorDB的一次Stan生成：10链，每链总迭代20000、预热10000、每10步保留一次，得到每链1000个相关样本；不是10000个独立实验。六个连续函数的MCSE为有限估计，不能约束共同偏差。事件样本全部为零，MCSE仍未定。

基于全局对数凹性支持平面的尾部积分不等式，256个角扇区给出的$R=12$尾质量/计算归一化常数为$7.47\times10^{-30}$。这是浮点计算的尾部上界估计，不包含区间算术认证、内部求积或舍入误差。六个连续均值从$R=6$到12的最大变化约$1.87\times10^{-9}$；事件由$4.34669\times10^{-11}$变为$4.76721\times10^{-11}$。完整12组值及各函数尾部估计保存在参考伴随数据中。

\subsection{L1/L2有限参考及已有偏移分析}
\begin{table}[htbp]\centering\small
\begin{tabular}{llrrrr}\toprule
目标 & 函数 & 参考均值 & 批均值MCSE & 拟合间MCSE & 采用MCSE\\\midrule
L1 & $\theta_1$ & $-0.278065574$ & $3.89\times10^{-4}$ & $2.59\times10^{-4}$ & $3.89\times10^{-4}$\\
L1 & $\theta_2$ & $-0.678895426$ & $5.00\times10^{-4}$ & $3.55\times10^{-4}$ & $5.00\times10^{-4}$\\
L1 & $\sigma(\theta_1)$ & $0.431794752$ & $9.41\times10^{-5}$ & $6.09\times10^{-5}$ & $9.41\times10^{-5}$\\
L1 & $\ind(\theta_1>0)$ & $0.109191895$ & $6.59\times10^{-4}$ & $7.64\times10^{-4}$ & $7.64\times10^{-4}$\\
L2 & $\theta_1$ & $-0.311176772$ & $8.89\times10^{-5}$ & $8.69\times10^{-5}$ & $8.89\times10^{-5}$\\
L2 & $\theta_2$ & $-0.600525763$ & $1.02\times10^{-4}$ & $7.79\times10^{-5}$ & $1.02\times10^{-4}$\\
L2 & $\sigma(\theta_1)$ & $0.422880789$ & $2.17\times10^{-5}$ & $2.12\times10^{-5}$ & $2.17\times10^{-5}$\\
L2 & $\ind(\theta_1>0)$ & $0$ & 未定 & 未定 & 未定\\
\bottomrule\end{tabular}
\caption{L1/L2逐函数有限参考。每模型四份独立四链拟合，每链保留16384步，共262144个相关样本；参考是归档reference-v1的既有BlackJAX NUTS结果。L2符号事件的原始机械方差计算为零，正式不确定性仍保留未定。}
\label{tab:reference-logistic}
\end{table}
四份参考拟合的种子与主要实验分离，模型数据和目标哈希经复用审查一致；每链预热1024步，最大树深10。四份拟合未记录发散，但这些检查不能排除参考的共同偏差。MCSE取批均值与拟合间估计的较大者。L1符号事件出现28624次，L2为零；后者不能作为精确零概率参考。

原正式分析已对有限MCMC参考作共同偏移$\mu^\star\pm2\,\mathrm{MCSE}$敏感性计算。在4096预算CPU顺序MALA与CPU Pyro NUTS的比较中，L1四函数及L2前三函数的损失差在该范围内均保持正号；但每格仅六个共同有效重复，未达到20份的区间计算下限。L2符号事件保持未定。逐函数上下限保存在参考伴随JSON；这一偏移范围不是置信区间，且没有扣除参考方差。

\subsection{W1比较对参考偏移的敏感性}
在相同可用重复集合和同一尺度$s_f=1$下，记两方法平均平方损失差为$D_f(c)$，参考从$c$移至$c+e$时有
\[
D_f(c+e)=D_f(c)-2e\left(\overline{\widehat\mu}_{A,f}-\overline{\widehat\mu}_{B,f}\right).
\]
这里$A$是CPU顺序MALA、$B$是CPU Pyro NUTS四进程工作流。本表固定共同集合以对应配对比较；改用各自集合会改变所比较的估计对象。若尺度不同，共同平方项一般不再抵消。下表固定4096预算和全部24份共同有效重复。$e_0$为使点损失差变为零的带符号参考偏移；若两方法平均估计相同，该阈值不存在。

\begin{table}[htbp]\centering\small
\begin{tabular}{lrrrr}\toprule
函数 & $D_f(c)$ & $e_0$ & 改用MCMC参考 & $\pm2$ MCSE偏移范围\\\midrule
$\alpha$ & $3.06\times10^{-7}$ & $8.55\times10^{-4}$ & $2.58\times10^{-7}$ & $-1.47\times10^{-7}$ 至 $7.58\times10^{-7}$\\
$\beta$ & $5.88\times10^{-7}$ & $-6.56\times10^{-4}$ & $1.34\times10^{-6}$ & $-1.27\times10^{-6}$ 至 $2.44\times10^{-6}$\\
$\alpha^2$ & $4.90\times10^{-7}$ & $0.00107$ & $4.28\times10^{-7}$ & $-2.11\times10^{-7}$ 至 $1.19\times10^{-6}$\\
$\beta^2$ & $1.02\times10^{-6}$ & $8.41\times10^{-4}$ & $2.35\times10^{-6}$ & $-2.27\times10^{-6}$ 至 $4.32\times10^{-6}$\\
$p(0)$ & $1.57\times10^{-8}$ & $1.94\times10^{-4}$ & $1.31\times10^{-8}$ & $-7.67\times10^{-9}$ 至 $3.90\times10^{-8}$\\
$p(100)$ & $1.67\times10^{-9}$ & $-1.24\times10^{-5}$ & $3.43\times10^{-8}$ & $-4.08\times10^{-8}$ 至 $4.41\times10^{-8}$\\
\bottomrule\end{tabular}
\caption{W1的参考敏感性。正号表示本格NUTS工作流的平均平方损失较小。末列只是以独立MCMC的MCSE作为假设偏移尺度，不是求积误差估计，也不是置信区间。}
\label{tab:reference-sensitivity}
\end{table}
六个连续函数在全部12个保存求积值以及独立MCMC参考点下保持相同点排序；在假设$\pm2$ MCSE偏移内，六者均可改变符号。因此数据支持对已检查参考点的稳定性，尚不提供对未知真值的认证排序。W1的$p(100)$即使在固定参考下，其原逐点BCa损失差区间也包含零。其余五个连续函数的固定参考区间不包含零；这些区间不传播求积参考的不确定性，也未校正多重比较。

所有W1拟合的$\ind(\beta>0)$均值都是零；图中的非零平方差完全由共同数值参考决定。改变共同参考仍使两方法该函数的点损失差为零，这既不确认极小尾概率的估计精度，也不解决常量函数诊断的未定状态。

\subsection{W1模型给出的后验解释}
\begin{table}[htbp]\centering\small
\begin{tabular}{lrrr}\toprule
参数或概率 & 后验均值 & 95\%等尾后验区间 & 均值MCSE\\\midrule
$\alpha$ & 0.606710 & [0.488359, 0.723356] & $6.33\times10^{-4}$\\
$\beta$ & -0.622143 & [-0.810761, -0.429960] & $0.00104$\\
$p(0)$ & 0.647069 & [0.619720, 0.673346] & $1.44\times10^{-4}$\\
$p(100)$ & 0.496146 & [0.465215, 0.526176] & $1.58\times10^{-4}$\\
\bottomrule\end{tabular}
\caption{独立既有参考样本的应用摘要。距离以100米为单位，$p(0)=\sigma(\alpha)$，$p(100)=\sigma(\alpha+\beta)$。区间为10000个保存相关样本的经验2.5\%与97.5\%分位数（R type 7等价规则）；它描述后验不确定性，均值MCSE描述有限计算误差。分位数本身的估计误差没有另行量化。}
\label{tab:wells-posterior-application}
\end{table}
该二维模型中，较远距离与较低的换井概率相联系。以上为已指定模型及其参考样本的后验描述，不把距离系数解释为因果效应；也没有从正式实验中挑选表现较好的拟合。

\noindent\textbf{来源与重建。}表格由随附的参考精度脚本从W1求积与参考档案、L1/L2参考摘要，以及Mac独立重建的正式统计生成。脚本校验目标、参考二进制和读取的统计文件哈希；参考伴随JSON同时保存完整逐函数数值、所有12个求积配置、两预算配对结果和来源SHA。没有新增MCMC或更换原正式参考。

